\chapter{2011年福建卷（文）}

\section{单选题}

\begin{problem}
若集合 \(M = \{-1, 0, 1\}\)，\(N = \{0, 1, 2\}\)，则 \(M \cap N\) 等于（\quad）

\choices
    {\(\{0, 1\}\)}
    {\( \{-1,0,1\} \)}
    {\(\{0,1,2\}\)}
    {\(\{-1,0,1,2\}\)}
\end{problem}

\begin{problem}
i 是虚数单位，\( 1 + \i^{2} \) 等于（\quad）

\choices
    {i}
    {-i}
    {\(1 + \i\)}
    {\(1 - \i\)}
\end{problem}

\begin{problem}
若 \(a \in \R\)，则 \(a^a = 1\) 是 \([a] = 1\) 的（\quad）

\choices
    {充分而不必要条件}
    {必要而不充分条件}
    {充要条件}
    {既不充分也不必要条件}
\end{problem}

\begin{problem}
某校选修乒乓球课程的学生中，高一年级有 30 名，高二年级有 40 名. 现用分层抽样的方法在这 70 名学生中抽取一个样本，已知在高一年级的学生中抽取了 6 名，则在高二年级的学生中应抽取的人数为（\quad）

\choices
    {6}
    {8}
    {10}
    {12}
\end{problem}

\begin{problem}
阅读如图所示的程序框图，运行相应的程序，输出的结果是（\quad）

\choices
    {3}
    {11}
    {38}
    {123}
\end{problem}

\begin{problem}
若关于 \(x\) 的方程 \(x^{2} + mx + 1 = 0\) 有两个不相等的实数根，则实数 \(m\) 的取值范围是（\quad）

\choices
    {\((-1, 1)\)}
    {\((-2, 2)\)}
    {\((- \infty, -2) \cup (2, +\infty)\)}
    {\((-\infty, -1) \cup (1, +\infty)\)}
\end{problem}

\begin{problem}
如图，矩形 \(ABCD\) 中，点 \(E\) 为边 \(CD\) 的中点，若在矩形 \(ABCD\) 内部随机取一个点 \(Q\)，则点 \(Q\) 取自 \(\triangle ABE\) 内部的概率等于（\quad）

\choices
    {\( \frac{1}{4} \)}
    {\( \frac{1}{3} \)}
    {\( \frac{1}{2} \)}
    {\( \frac{2}{3} \)}
\end{problem}

\begin{problem}
已知函数 \( f(x) = \begin{cases} 2^x, & x > 0, \\ x + 1, & x \leqslant 0. \end{cases} \) 若 \( f(a) + f(1) = 0 \)，则实数 \( a \) 的值等于（\quad）

\choices
    {\(-3\)}
    {\(-1\)}
    {1}
    {3}
\end{problem}

\begin{problem}
若 \(\alpha \in \left(0, \frac{\pi}{2}\right)\)，且 \(\sin^2\alpha + \cos 2\alpha = \frac{1}{4}\)，则 \(\tan \alpha\) 的值等于（\quad）

\choices
    {\(\frac{\sqrt{2}}{2}\)}
    {\( \frac{\sqrt{3}}{3} \)}
    {\(\sqrt{2}\)}
    {\(\sqrt{3}\)}
\end{problem}

\begin{problem}
若 \(a > 0, b > 0\)，且函数 \(f(x) = 4x^3 - ax^2 - 2bx + 2\) 在 \(x = 1\) 处有极值，则 \(ab\) 的最大值等于（\quad）

\choices
    {2}
    {3}
    {6}
    {9}
\end{problem}

\begin{problem}
设圆锥曲线 \(T\) 的两个焦点分别为 \(F_{1}, F_{2}\). 若曲线 \(T\) 上存在点 \(P\) 满足 \(|PF_{1}|: |F_{1}F_{2}|: |PF_{2}| = 4:3:2\)，则曲线 \(T\) 的离心率等于（\quad）

\choices
    {\(\frac{1}{2}\) 或 \(\frac{3}{2}\)}
    {\(\frac{2}{3}\) 或 2}
    {\(\frac{1}{2}\) 或 2}
    {\(\frac{2}{3}\) 或 \(\frac{3}{2}\)}
\end{problem}

\begin{problem}
在整数集 \(\Z\) 中，被 5 除所得余数为 \(\mathbf{k}\) 的所有整数组成一个“夹”，记为 \([k]\)，即 \([k] = \{5n + k \mid n \in \Z\}, k = 0, 1, 2, 3, 4\). 给出如下四个结论：\circled{1} \(2011 \in [1]\)；\circled{2} \(-3 \in [2]\)；\circled{3} \(\Z = [0] \cup [1] \cup [2] \cup [3] \cup [4]\)；\circled{4} “整数 \(a, b\) 属于同一类” 的充要条件是 “\(a - b \in [0]\)”. 其中，正确结论的个数是（\quad）

\choices
    {1}
    {2}
    {3}
    {4}
\end{problem}

\section{填空题}

\begin{problem}
若向量 \( a=(1,1) \) ，\( b=(-1,2) \) ，则 \( a \cdot b \) 等于 \fillinblank{} \fillinblank{}.
\end{problem}

\begin{problem}
若 \(\triangle ABC\) 的面积为 \(\sqrt{3}\)，\(BC = 2\)，\(C = 60^{\circ}\)，则边 \(AB\) 的长度等于 \fillinblank{} \fillinblank{}.
\end{problem}

\begin{problem}
如图，正方体 \(ABCD - A_{1}B_{1}C_{1}D_{1}\) 中，\(AB = 2\)，点 \(E\) 为 \(AD\) 的中点，点 \(F\) 在 \(CD\) 上，若 \(EF //\) 平面 \(A_{1}B_{1}C_{1}\) 则线段 \(EF\) 的长度等于 \fillinblank{} \fillinblank{}.
\end{problem}

\begin{problem}
商家通常依据“乐玩系数准则”确定商品销售价格，即根据商品的最低销售限价 a，最高销售限价 b (b > a) 以及实数 x (0 < x < 1) 确定实际销售价格 \( c = a + x (b - a) \) . 这里，x 被称为乐玩系数. 经验表明，最佳乐玩系数 x 恰好使得 \( (c - a) \) 是 \( (b - c) \) 和 \( (b - a) \) 的等比中项. 据此可得，最佳乐玩系数 x 的值等于 \fillinblank{} \fillinblank{}.
\end{problem}

\section{解答题}

\begin{problem}
已知等差数列 \(\{a_{n}\}\) 中，\(a_{1} = 1, a_{3} = -3\). (1) 求数列 \( \{a_{n}\} \) 的通项公式; (2) 若数列 \( \{a_{n}\} \) 的前 k 项和 \( S_{k} = -35 \) ，求 k 的值.
\end{problem}

\begin{problem}
如图，直线 \(l: y = x + b\) 与抛物线 \(C: x^2 = 4y\) 相切于点 \(A\). (1) 求实数 b 的值; (2) 求以点 A 为圆心，且与抛物线 C 的准线相切的圆的方程.
\end{problem}

\begin{problem}
某日用品按行业质量标准分成五个等级. 等级系数 \(X\) 依次为 1, 2, 3, 4, 5. 现从一批该日用品中随机抽取 20 件，对其等级系数进行统计分析，得到频率分布表如下: X 1 2 3 4 5 f a 0.2 0.45 b c (1) 若所抽取的 20 件日用品中，等级系数为 4 的恰有 3 件，等级系数为 5 的恰有 2 件，求 a, b, c 的值; (2) 在 (1) 的条件下，将等级系数为 4 的 3 件日用品记为 \(x_{1}, x_{2}, x_{3}\)，等级系数为 5 的 2 件日用品记为 \(y_{1}, y_{2}\). 现从 \(x_{1}, x_{2}, x_{3}, y_{1}, y_{2}\) 这 5 件日用品中任取两件 (假定每件日用品被取出的可能性相同)，写出所有可能的结果，并求这两件日用品的等级系数恰好相等的概率.
\end{problem}

\begin{problem}
设函数 \( f(\theta) = \sqrt{3}\sin \theta + \cos \theta \)，其中，角 \( \theta \) 的顶点与坐标原点重合，始边与 \( x \) 轴非负半轴重合，终边经过点 \( P(x, y) \)，且 \( 0 \leqslant \theta \leqslant \pi \). (1) 若点 \(P\) 的坐标为 \(\left(\frac{1}{2}, \frac{\sqrt{3}}{2}\right)\). 求 \(f(\theta)\) 的值; (2) 若点 \(P(x, y)\) 为平面区域 \(\Omega: \left\{ \begin{array}{l} x + y \geqslant 1, \\ x \leqslant 1, \\ y \leqslant 1 \end{array} \right.\) 上的一个动点，试确定角 \(\theta\) 的取值范围，并求函数 \(f(\theta)\) 的最小值和最大值.
\end{problem}

\begin{problem}
已知 \(a, b\) 为常数，且 \(a \neq 0\). 函数 \(f(x) = -ax + b + ax \ln x, f(\e) = 2\) (\(\e = 2.71828 \cdots\) 是自然对数的底数). (1) 求实数 \(b\) 的值; (2) 求函数 \( f(x) \) 的单调区间; (3) 当 \(a = 1\) 时，是否同时存在实数 \(m\) 和 \(M (m < M)\)，使得对每一个 \(t \in [m, M]\)，直线 \(y = t\) 与曲线 \(y = f(x) \left( x \in \left[ \frac{1}{2}, \e \right] \right)\) 都有公共点\(\perp\) 若存在. 求出最小的实数 \(m\) 和最大的实数 \(M\); 若不存在，说明理由.
\end{problem}

\begin{problem}
如图，四棱锥 \(P - ABCD\) 中，\(PA \perp\) 底面 \(ABCD\)，\(AB \perp AD\)，点 \(E\) 在线段 \(AD\) 上，且 \(CE // AB\). (1) 求证: \( CE \perp \) 平面 PAD; (2) 若 \(PA = AB = 1, AD = 3, CD = \sqrt{2}, \angle CDA = 45^{\circ}\)，求四棱锥 \(P - ABCD\) 的体积.
\end{problem}

